% =====================================================================
%  ThmSmith Stage 0 derivation skeleton.
%  Written automatically by the ThmSmith TS pipeline before Stage 0 runs.
%  Locked parameters: qid=pid\_continuous\_iv\_no\_monotonicity, specialization=v1
%
%  HOW TO USE THIS FILE
%  --------------------
%  - The preamble (everything above the `% --- BEGIN CONTENT ---` marker)
%    and the `\section{...}` headers below are fixed by the pipeline.
%    Do NOT rewrite or rename them; downstream stages (0.5 reviewer,
%    Stage 1, Stage 4d) grep for these section titles and macros.
%  - Each section contains exactly one `% TODO(stage0 ...): ...`
%    placeholder. Replace each TODO with the content described for that
%    section in `ThmSmith/tools/src/prompts/stage0_outputFormat.txt`.
%  - The `\paragraph{Locked parameters.}` block inside Formal setup must
%    pin the convention precisely. For Panel anchors mapped to the Q1
%    pool-mode, fill it VERBATIM from the Q1 block of
%    `ThmSmith/doc/panel_formula_question_generator_note_with_common_prompts.tex`
%    (Q2..Q6 templates have been retired). For free-form panel / ExactID /
%    PartialID anchors (the default), reproduce the locked-parameter block
%    from the Stage -1 proposal verbatim. Do not rephrase locked parameters;
%    Stage 1 quotes them.
%  - Removing a section header, renaming a macro, or rewriting the
%    preamble is a regression: the file must still match this skeleton
%    after you finish.
% =====================================================================

\documentclass[11pt]{article}

\usepackage{amsmath,amssymb,amsthm,mathtools}
\usepackage{enumitem}
\usepackage[colorlinks=true,linkcolor=blue,citecolor=blue,urlcolor=blue]{hyperref}

\newcommand{\E}{\mathbb{E}}
\renewcommand{\Pr}{\mathbb{P}}
\newcommand{\one}{\mathbf{1}}
\newcommand{\transpose}{\mathsf{T}}
\newcommand{\calH}{\mathcal{H}}
\newcommand{\calF}{\mathcal{F}}
\newcommand{\calR}{\mathcal{R}}
\newcommand{\R}{\mathbb{R}}
\newcommand{\plim}{\operatorname*{plim}}

\theoremstyle{definition}
\newtheorem{definition}{Definition}
\newtheorem{assumption}{Assumption}
\newtheorem{example}{Example}

\theoremstyle{plain}
\newtheorem{theorem}{Theorem}
\newtheorem{proposition}{Proposition}
\newtheorem{lemma}{Lemma}
\newtheorem{corollary}{Corollary}

\theoremstyle{remark}
\newtheorem{remark}{Remark}

\title{ThmSmith Stage 0 derivation: \texttt{pid\_continuous\_iv\_no\_monotonicity} / \texttt{v1}%
  \\ \large\textit{Continuous-IV bounds with calibrated local nonmonotonicity}}
\author{ThmSmith Stage 0 solver}
\date{\today}

\begin{document}
\maketitle

% --- BEGIN CONTENT ---  (everything above is fixed by the pipeline)

\begin{abstract}
This note studies the partial-identification functional obtained by projecting
globally path-compatible continuous-IV response laws onto the positive-responder
ratio
\[
\theta=\frac{\int_0^1 a_+(z)\,dz}{\int_0^1 c_+(z)\,dz}.
\]
The observable law is \(P\) for \((Y,D,Z)\), with \(Y\in[0,1]\), binary
\(D\), and scalar \(Z\in[0,1]\); the maintained restrictions are IV
validity, exclusion, differentiable reduced forms, full level matching to
\(P\), and a known envelope \(0\le c_-(z)\le b(z)\) on negative response
density.  The bounding functional is
\((\theta_L(P,b),\theta_U(P,b))\), the infimum and supremum of the same
ratio over all full-\(P\)-compatible response paths satisfying the envelope.
The per-conjecture stages will verify whether local signed-derivative
constraints paste sharply into one global response path and whether the usual
quasi-monotone integrated IV formula is robustly separated from the resulting
positive-responder identified set.
\end{abstract}

\section{Introduction}
Continuous instruments are often used because infinitesimal changes in the
instrument can reveal marginal treatment effects rather than a single
finite-instrument local average.  The standard interpretation requires a
monotone first stage: as the instrument increases, units may enter treatment
but may not exit.  In applications based on prices, distances, eligibility
scores, or policy generosity, this monotonicity condition is economically
strong; small changes in the instrument can plausibly move some units in the
opposite direction.

The problem addressed here is the following calibrated sensitivity question.
The observable derivatives
\[
p'(z)=\frac{d}{dz}\E[D\mid Z=z],
\qquad
m'(z)=\frac{d}{dz}\E[Y\mid Z=z]
\]
are signed mixtures of positive and negative marginal response types.  If the
negative response density is known only to be bounded by \(b(z)\), what can be
said about the average treatment effect among positive marginal responders?
The answer cannot be just the derivative ratio \(m'(z)/p'(z)\), because a
negative response margin contributes with the opposite sign to both derivatives.

This note makes three contributions.  First, it derives the exact local
signed-response interval implied by bounded outcomes and the envelope
\(0\le c_-\le b\).  Second, it defines the global bounding functional as an
optimization over full response paths that match the whole observable law
\(P_{Y,D\mid Z=z}\), not merely the derivative pair \((p',m')\).  Third, it
separates what is proved from what remains a genuine sharpness issue:
local derivative constraints give an outer relaxation, while global
path-compatibility and endpoint attainment determine whether the closed
interval is fully sharp at its endpoints.  The result is useful for formula
design because it identifies the exact objects that a continuous-IV
monotonicity-violation bound must optimize and the conditions under which the
usual local-IV formula is recovered.

\paragraph{Motivating conjecture(s).}
\paragraph{Conjecture 1.}\label{conj:sharp-pasting}
Under Assumptions~\ref{ass:continuous-iv}--\ref{ass:path-compatible}, the sharp identified set \(\Theta_{\mathrm{CIV}}(P,b)\) is exactly the interval obtained by optimizing
\[
\frac{\int_0^1 a_+(z)\,dz}{\int_0^1 (p'(z)+f(z))\,dz}
\]
over measurable functions \(f\) and \(a_+\) satisfying the pointwise constraints in Theorem~\ref{thm:local-envelope}, with \(0\le f\le b\), together with the following level constraints: \(p(z)=p(0)+\int_0^z p'(u)\,du\), \(m(z)=m(0)+\int_0^z m'(u)\,du\), the induced conditional law of \((Y(D(z)),D(z))\) equals \(P_{Y,D\mid Z=z}\) for every \(z\), the signed first-stage derivative measure is \(p'(z)\,dz\), and the total-variation response measure \((c_++c_-)\,dz=(p'+2f)\,dz\) is generated by one bounded-variation binary response path.  Moreover, the endpoints are attained by extreme full-data signed-response laws, so the interval is sharp rather than an outer relaxation.

\paragraph{Conjecture 2.}\label{conj:strict-separation}
There exists a nonempty open set of observable laws \(P\), baseline envelope directions \(b_0\), and constants \(\delta>0\) and \(\eta>0\) satisfying Assumptions~\ref{ass:continuous-iv}--\ref{ass:path-compatible} with \(\int_0^1 p'(z)\,dz>0\) such that the quasi-monotone integrated formula
\[
q_{\mathrm{IV}}(P)=\frac{\int_0^1 m'(z)\,dz}{\int_0^1 p'(z)\,dz}
\]
has distance at least \(\delta\) from \(\Theta_{\mathrm{CIV}}(P,b_0)\).  This separation is locally robust: for every observable law \(\widetilde P\) and envelope direction \(\widetilde b_0\) in an \(L^1\)-neighborhood of radius \(\eta\), \(\operatorname{dist}(q_{\mathrm{IV}}(\widetilde P),\Theta_{\mathrm{CIV}}(\widetilde P,\widetilde b_0))\ge\delta/2\), and the robustness curve \(t\mapsto\mathcal R_{\widetilde P}(t)\) has a positive, locally Lipschitz breakdown radius \(\tau^\star(\widetilde P;\widetilde b_0)>1\).  A target witness class partitions \([0,1]\) into two positive-measure regions \(A\) and \(B\), has \(p'\ge\kappa>0\) on both regions, places \(b_0>0\) only on \(B\), and assigns positive outcome variation so that positive responders on \(A\) and bounded negative responders on \(B\) have separated treatment-effect numerators with all inequalities in Theorem~\ref{thm:local-envelope} slack.  The separation is generated by negative response density concentrated where defier treatment effects differ from positive-responder treatment effects, so the signed derivative ratio is not a member of the positive-responder sharp set even though the observed first stage remains positive.

\section{Related work and positioning}
The classical binary-IV reference is \cite{ImbensAngrist1994}, which identifies
a local average treatment effect under independence, exclusion, relevance, and
monotonicity.  This note keeps the positive-responder causal target but replaces
exact monotonicity by a known bound on the local negative response density.
\cite{Vytlacil2002} shows that the LATE monotonicity condition is equivalent,
under the usual IV restrictions, to a latent-index threshold-crossing model;
here \(c_-(z)\) is precisely the local object that records departures from that
threshold-crossing structure.  \cite{HeckmanVytlacil2005} derive the local
instrumental-variable and marginal-treatment-effect formula using derivatives
of the reduced forms; Theorem~\ref{thm:zero-envelope} below verifies that
the proposed bound collapses to that formula when \(b=0\).

\cite{Klein2010} studies instrumental variables without monotonicity and shows
how local violations create bias terms in monotonicity-based estimands.  The
current note differs by treating the negative response density as a calibrated
partial-identification restriction and by deriving a feasible-set projection
rather than a bias expression around a point estimand.  \cite{KennedyLorchSmall2019}
develop robust estimation of the local-IV curve for continuous instruments; the
object here is not an estimator of that curve but a population identified set
when the curve's monotone-complier interpretation is relaxed.  \cite{Marx2024}
gives sharp bounds for marginal treatment effects and linear functionals inside
a monotone latent-index selection model; the difference is that the present
response path may have both positive and negative variation, with the negative
part bounded by \(b\).

The partial-identification geometry is closest to Manski worst-case completion
and finite instrumental-variable linear programs.  \cite{Manski1990} gives the
canonical bounded-outcome logic: when the missing counterfactual component is
not restricted, the identified set is obtained by assigning it its extreme
support values.  \cite{BalkePearl1997} derive sharp binary-IV bounds through a
finite latent-type linear program; the continuous-IV program here is its
infinite-dimensional signed-response analogue with a path-compatibility
constraint across adjacent instrument values.  \cite{Kitagawa2015} characterizes
testable implications of IV validity and monotonicity, which is relevant because
the feasibility of the present derivative program is also a compatibility
condition between observables and latent response restrictions.

Several monotonicity-relaxation papers are adjacent but do not solve this
problem.  \cite{DeChaisemartin2017} obtains LATE-like interpretation without
full monotonicity by imposing structure connecting compliers and defiers, while
this note imposes a quantitative envelope and otherwise allows arbitrary
bounded treatment-effect heterogeneity.  \cite{Yap2022} studies sensitivity of
policy-relevant parameters to monotonicity violations with finite-instrument
defier-to-complier calibration; the present envelope is indexed continuously by
the instrument margin.  \cite{MogstadSantosTorgovitsky2018} use IV-like moments
and MTE restrictions to bound policy-relevant treatment parameters; the current
\(\Theta_{\mathrm{CIV}}(P,b)\) is a candidate robustness input for that policy
evaluation agenda.  The in-repository bounded-defier and ordinal-IV proposal
families are relevant because they exposed the same distinction between local
cell accounting and global latent-response compatibility; the present note
keeps that distinction explicit for a scalar continuous instrument.

\section{Formal setup}
Let \(P\) be the observable law of \((Y,D,Z)\), where \(Y\in[0,1]\),
\(D\in\{0,1\}\), and \(Z\in[0,1]\).  The instrument has density \(g\) with
\(0<\underline g\le g(z)\le \overline g<\infty\), so conditioning on \(Z=z\)
is interpreted through a regular conditional distribution \(P_{Y,D\mid Z=z}\).
Write
\[
p(z)=\E_P[D\mid Z=z],
\qquad
m(z)=\E_P[Y\mid Z=z].
\]
The observable map is therefore
\[
(Y,D,Z)\mapsto P,\qquad
P\mapsto \{P_{Y,D\mid Z=z}:z\in[0,1]\},\quad p,\quad m.
\]

A full-data law consists of a binary response path \(z\mapsto D(z)\) of bounded
variation and potential outcomes \(Y(1),Y(0)\in[0,1]\).  Exclusion means that
the observed outcome at instrument value \(z\) is
\[
Y(D(z))=Y(1)D(z)+Y(0)(1-D(z)).
\]
The positive and negative variation measures of the response path, aggregated
over the population, are assumed absolutely continuous with densities
\(c_+\) and \(c_-\).  Their treatment-effect numerator measures are also
absolutely continuous with densities \(a_+\) and \(a_-\).  Thus the derivative
identities are
\[
p'(z)=c_+(z)-c_-(z),
\qquad
m'(z)=a_+(z)-a_-(z)
\]
for almost every \(z\).

For a fixed envelope \(b\in L^1_+([0,1])\), let \(\mathfrak L(P,b)\) be the set
of all full-data laws satisfying the assumptions in Section~\ref{sec-assumptions}
below, matching \(P_{Y,D\mid Z=z}\) for every \(z\), and satisfying
\(0\le c_-(z)\le b(z)\) almost everywhere.  The identified set is
\[
\Theta_I(P,b)
=
\left\{
\frac{\int_0^1 a_+(z)\,dz}{\int_0^1 c_+(z)\,dz}
:
L\in\mathfrak L(P,b),\ \int_0^1 c_+(z)\,dz>0
\right\}.
\]
The notation \(\Theta_{\mathrm{CIV}}(P,b)\) denotes the same set, emphasizing
that it is the continuous-IV sharp projection over full-\(P\)-compatible
response paths.

\paragraph{Locked parameters.}
Free-form partial-identification anchor.  The locked tuple is
\[
(\texttt{qid},\texttt{specialization},P,\theta,\text{identifying restrictions},\text{bounding functional}).
\]
Here \(\texttt{qid}=\texttt{pid\_continuous\_iv\_no\_monotonicity}\) and \(\texttt{specialization}=\texttt{v1}\).  \(P\) is the observable law of \((Y,D,Z)\), where \(Y\in[0,1]\), \(D\in\{0,1\}\), and the scalar instrument \(Z\) has support \(\mathcal Z=[0,1]\) with a positive density; write \(p(z)=\E[D\mid Z=z]\) and \(m(z)=\E[Y\mid Z=z]\).  The parameter \(\theta\) is the average treatment effect for infinitesimal positive responders induced along the \(Z\)-margin,
\[
\theta=\frac{\int_{\mathcal Z} a_+(z)\,dz}{\int_{\mathcal Z} c_+(z)\,dz},
\]
where \(c_+(z)\) is the positive first-stage response density and \(a_+(z)\) is the corresponding signed treatment-effect numerator density.  The identifying restrictions are instrument independence and exclusion, bounded potential outcomes \(Y(d)\in[0,1]\) so numerator densities may be signed but obey \(-c_\pm(z)\le a_\pm(z)\le c_\pm(z)\), differentiability of the observable reduced forms, level matching to the full observable law \(P\), and a calibrated nonmonotonicity envelope \(0\le c_-(z)\le b(z)\) on the negative first-stage response density with \(p'(z)=c_+(z)-c_-(z)\).  The bounding functional is the sharp projection of the signed-response feasible set onto \(\theta\), denoted \(\Theta_{\mathrm{CIV}}(P,b)\).

\section{Assumptions}\label{sec-assumptions}
\begin{assumption}[Observable support and bounded outcomes]\label{ass:continuous-iv}
\(Z\) has support \([0,1]\) and a density bounded away from zero and infinity.
\(D\in\{0,1\}\), \(Y(0),Y(1)\in[0,1]\) almost surely, and hence
\(Y\in[0,1]\).
\end{assumption}

\begin{assumption}[Instrument independence and exclusion]\label{ass:iv-valid}
\(Z\) is independent of the full latent response path
\(\{D(z):z\in[0,1]\}\) and the potential outcomes \((Y(1),Y(0))\).
The observed variables obey \(D=D(Z)\) and \(Y=Y(D(Z))\).
\end{assumption}

\begin{assumption}[Differentiable reduced forms]\label{ass:differentiable}
The reduced forms \(p(z)=\E[D\mid Z=z]\) and
\(m(z)=\E[Y\mid Z=z]\) are absolutely continuous on \([0,1]\), with
derivatives \(p',m'\in L^1([0,1])\).  Endpoint statements use the
corresponding one-sided absolutely continuous representatives.
\end{assumption}

\begin{assumption}[Signed response decomposition]\label{ass:signed-response}
The population positive and negative variation measures of the response path
are absolutely continuous with nonnegative densities
\(c_+,c_-\in L^1([0,1])\).  The corresponding treatment-effect numerator
measures have densities \(a_+,a_-\in L^1([0,1])\) satisfying, almost
everywhere,
\[
p'=c_+-c_-,
\qquad
m'=a_+-a_-,
\qquad
-c_+\le a_+\le c_+,
\qquad
-c_-\le a_-\le c_-.
\]
\end{assumption}

\begin{assumption}[Calibrated nonmonotonicity envelope]\label{ass:envelope}
The negative response density is bounded by a known envelope
\(b\in L^1_+([0,1])\):
\[
0\le c_-(z)\le b(z)
\quad\text{for almost every }z.
\]
\end{assumption}

\begin{assumption}[Full-\(P\) level matching and global path compatibility]\label{ass:path-compatible}
The densities \(c_+,c_-\) and \(a_+,a_-\) are induced by one measurable
bounded-variation binary response path \(z\mapsto D(z)\) and one pair of
potential outcomes \((Y(1),Y(0))\).  The same full-data law reproduces the
observable conditional law \(P_{Y,D\mid Z=z}\) for every \(z\), not only the
two derivative functions \(p'\) and \(m'\).
\end{assumption}

\begin{assumption}[Positive positive-responder mass]\label{ass:positive-mass}
Every feasible full-data law considered in \(\mathfrak L(P,b)\) satisfies
\[
\int_0^1 c_+(z)\,dz>0.
\]
For the zero-envelope local-IV reduction, assume additionally that
\(p'(z)\ge \kappa>0\) almost everywhere for some \(\kappa>0\).
\end{assumption}

\section{Estimand-defining functional}
For \(P\) and \(b\), define the global signed-response shadow set
\(\mathfrak S(P,b)\) as the set of all pairs \((f,u)\in L^1_+([0,1])\times
L^1([0,1])\) for which there exists a full-data law in \(\mathfrak L(P,b)\)
such that
\[
f(z)=c_-(z),
\qquad
u(z)=a_+(z)
\quad\text{almost everywhere.}
\]
For any \((f,u)\in\mathfrak S(P,b)\), the positive response density is
\[
c_+(z)=p'(z)+f(z),
\]
and the numerator ratio is
\[
R_P(f,u)=
\frac{\int_0^1 u(z)\,dz}
     {\int_0^1 (p'(z)+f(z))\,dz}.
\]
Assumption~\ref{ass:positive-mass} makes the denominator strictly positive on
\(\mathfrak S(P,b)\).

The bounding functional is the closed interval
\[
\theta_L(P,b)=\inf_{(f,u)\in\mathfrak S(P,b)} R_P(f,u),
\qquad
\theta_U(P,b)=\sup_{(f,u)\in\mathfrak S(P,b)} R_P(f,u),
\]
with the convention that the interval is empty when \(\mathfrak S(P,b)\) is
empty.  This is a sharp closed-interval projection of the globally compatible
feasible set.  It should be distinguished from the local derivative relaxation
\(\mathfrak S_{\mathrm{loc}}(P,b)\), which keeps only the pointwise inequalities
derived in Theorem~\ref{thm:local-envelope} and drops full-\(P\) matching
and path compatibility.  Optimizing over \(\mathfrak S_{\mathrm{loc}}(P,b)\)
gives an outer bound; it need not be sharp for \(\Theta_I(P,b)\).

When \(\int_0^1 p'(z)\,dz>0\), define the quasi-monotone integrated IV formula
\[
q_{\mathrm{IV}}(P)=
\frac{\int_0^1 m'(z)\,dz}{\int_0^1 p'(z)\,dz}.
\]
This is an observable functional of \(P\), but it is not the target and is not
generally an element of the positive-responder identified set.

\section{Target causal estimand}
The causal target is the average treatment effect for infinitesimal positive
responders induced along the instrument margin:
\[
\theta(L)=
\frac{\int_0^1 a_+(z)\,dz}
     {\int_0^1 c_+(z)\,dz}.
\]
Here \(c_+(z)\,dz\) is the population measure of units whose treatment path
switches from \(0\) to \(1\) at margins in \(dz\), and \(a_+(z)\,dz\) is the
same measure weighted by the unit-level treatment effect \(Y(1)-Y(0)\).  The
target is a scalar attached to a full-data law \(L\), not to the observable law
\(P\) alone.  Partial identification asks which values of \(\theta(L)\) are
compatible with \(P\) and the maintained restrictions; the answer is a set or
interval, not an equality between \(\theta\) and an observable derivative ratio.

\section{Main theorems}
\begin{theorem}[Conjecture 1: sharp global pasting of local envelopes]\label{thm:sharp-pasting}
Under Assumptions~\ref{ass:continuous-iv}--\ref{ass:path-compatible}, the sharp identified set \(\Theta_{\mathrm{CIV}}(P,b)\) is exactly the interval obtained by optimizing
\[
\frac{\int_0^1 a_+(z)\,dz}{\int_0^1 (p'(z)+f(z))\,dz}
\]
over measurable functions \(f\) and \(a_+\) satisfying the pointwise constraints in Theorem~\ref{thm:local-envelope}, with \(0\le f\le b\), together with the following level constraints: \(p(z)=p(0)+\int_0^z p'(u)\,du\), \(m(z)=m(0)+\int_0^z m'(u)\,du\), the induced conditional law of \((Y(D(z)),D(z))\) equals \(P_{Y,D\mid Z=z}\) for every \(z\), the signed first-stage derivative measure is \(p'(z)\,dz\), and the total-variation response measure \((c_++c_-)\,dz=(p'+2f)\,dz\) is generated by one bounded-variation binary response path.  Moreover, the endpoints are attained by extreme full-data signed-response laws, so the interval is sharp rather than an outer relaxation.

Equivalently, with \(\mathfrak S(P,b)\), \(R_P\), \(\theta_L(P,b)\), and
\(\theta_U(P,b)\) defined in Section~7 on the positive-mass domain,
\[
\Theta_{\mathrm{CIV}}(P,b)
=
\bigl[\theta_L(P,b),\theta_U(P,b)\bigr],
\qquad
\theta_L(P,b)=\min_{(f,u)\in\mathfrak S(P,b)}R_P(f,u),
\quad
\theta_U(P,b)=\max_{(f,u)\in\mathfrak S(P,b)}R_P(f,u).
\]
The endpoints are also the limits of the finite mesh linear-fractional
programs whose variables are latent binary response paths on a partition of
\([0,1]\), binned potential-outcome types in \([0,1]^2\), and nonnegative
type weights matching the observed binned laws \(P_{Y,D\mid Z=z_j}\), the
cell first-stage increments, the cell outcome increments, and the integrated
envelope constraints.
\end{theorem}

The theorem says that the local envelope in Theorem~\ref{thm:local-envelope}
is not merely a necessary derivative accounting identity.  Once the candidate
densities are required to come from one full-data bounded-variation response
path that matches the entire conditional law \(P_{Y,D\mid Z=z}\), optimizing
the positive-responder numerator ratio over those candidates is exactly the
sharp partial-identification problem.

\begin{theorem}[Conjecture 2: strict structural separation]\label{thm:strict-separation}
Conjecture 2 is confirmed: the two-region tagged-response witness class below gives a nonempty \(L^1\)-open neighborhood on which \(q_{\mathrm{IV}}\) is separated from \(\Theta_{\mathrm{CIV}}(P,b_0)\) and has locally Lipschitz breakdown radius strictly larger than one.

At the center of the witness class, let
\[
A=[0,1/2],\qquad B=(1/2,1],
\]
and let the response densities and numerator densities be constant on these two regions as follows:
\[
\begin{array}{c|c|c|c|c}
\text{region/type} & c_+ & a_+ & c_- & a_-\\
\hline
A\text{ visible positive responders} & 1/10 & 2/25 & 0 & 0\\
B\text{ visible positive responders} & 1/10 & 2/25 & 0 & 0\\
B\text{ visible negative responders} & 0 & 0 & 1/25 & 9/250
\end{array}
\]
The baseline envelope is
\[
b_0(z)=0\quad(z\in A),\qquad b_0(z)=21/500\quad(z\in B).
\]
The observable derivatives are therefore
\[
p'(z)=1/10,\quad m'(z)=2/25\quad(z\in A),
\]
and
\[
p'(z)=3/50,\quad m'(z)=11/250\quad(z\in B).
\]
The witness law also contains, only on \(B\), an unobserved zero-effect cancelling response component of density \(s\ge0\): equal masses of \(0\to1\) and \(1\to0\) movers with \(Y(1)=Y(0)=1/2\).  This component leaves \(P_{Y,D\mid Z=z}\), \(p'\), and \(m'\) unchanged, contributes \(s\) to both \(c_+\) and \(c_-\), and contributes zero to both \(a_+\) and \(a_-\).  Full-\(P\) level matching and the separated outcome tags force the visible entries in the table and leave \(s\) as the only free scalar.  Under the baseline envelope,
\[
0\le s\le b_0|_B-1/25=1/500.
\]
Consequently
\[
\Theta_{\mathrm{CIV}}(P^\star,b_0^\star)
=
\left[
\frac{2/25}{1/10+(1/2)(1/500)},\,
\frac{2/25}{1/10}
\right]
=
\left[\frac{80}{101},\,\frac45\right],
\]
whereas
\[
q_{\mathrm{IV}}(P^\star)
=
\frac{(1/2)(2/25)+(1/2)(11/250)}
     {(1/2)(1/10)+(1/2)(3/50)}
=\frac{31}{40}.
\]
Thus
\[
\operatorname{dist}\!\left(q_{\mathrm{IV}}(P^\star),
\Theta_{\mathrm{CIV}}(P^\star,b_0^\star)\right)
=
\frac{80}{101}-\frac{31}{40}
=\frac{69}{4040}
>\frac1{100}.
\]
For the robustness curve \(\mathcal R_{P^\star}(t)=\Theta_{\mathrm{CIV}}(P^\star,tb_0^\star)\), the available cancelling density on \(B\) is
\[
0\le s\le t(21/500)-1/25.
\]
The value \(q_{\mathrm{IV}}(P^\star)=31/40\) is reached exactly when
\[
\frac{2/25}{1/10+s/2}=\frac{31}{40},
\qquad\text{namely}\qquad s=\frac1{155}.
\]
Hence
\[
\tau^\star(P^\star;b_0^\star)
=
\frac{1/25+1/155}{21/500}
=\frac{240}{217}>1.
\]
The same inequalities continue to hold throughout a sufficiently small
\(L^1\)-open neighborhood of the tagged cell-flow vector and envelope direction.
For that neighborhood one may take \(\delta=1/100\) and some \(\eta>0\) small
enough that the perturbed lower endpoint, \(q_{\mathrm{IV}}\), and
\(\tau^\star\) move by less than one fourth of the displayed strict margins.
\end{theorem}

The example isolates the reason the signed derivative ratio is not the
positive-responder target.  On both \(A\) and \(B\), positive responders have
treatment-effect numerator \(a_+/c_+=4/5\); on \(B\), however, the observed
outcome derivative subtracts the negative-responder numerator
\(a_-=9/250\), while the observed first-stage derivative subtracts the
negative-responder density \(c_-=1/25\).

\section{Interpretation}
\paragraph{Interpretation for Conjecture~\ref{conj:sharp-pasting}.}
The theorem says that the local envelope in Theorem~\ref{thm:local-envelope}
is not merely a necessary derivative accounting identity.  Once the candidate
densities are required to come from one full-data bounded-variation response
path that matches the entire conditional law \(P_{Y,D\mid Z=z}\), optimizing
the positive-responder numerator ratio over those candidates is exactly the
sharp partial-identification problem.  The causal estimand remains
\[
\theta(L)=
\frac{\int_0^1 a_+(z)\,dz}{\int_0^1 c_+(z)\,dz}
\]
for a full-data law \(L\); the observable law \(P\) identifies only the set of
values generated by all globally compatible laws.  The pointwise inequalities
are useful because they reduce every feasible law to the two observable
derivatives \(p'\), \(m'\), the unknown negative response density \(f=c_-\),
and the positive-responder numerator \(u=a_+\).  The path and level
constraints are equally essential: dropping them would give only a local outer
relaxation, not the sharp set.

The endpoint formula is computable from observed-population inputs in the
following population sense.  For a partition
\(\pi=\{0=z_0<z_1<\cdots<z_J=1\}\) and an outcome grid
\(\mathcal H=\{h_1,\ldots,h_H\}\), solve the finite linear-fractional program
with variables
\[
\lambda_{r,h_0,h_1}\ge0,
\qquad
r\in\{0,1\}^{J+1},\quad (h_0,h_1)\in\mathcal H^2,
\]
where \(\sum_{r,h_0,h_1}\lambda_{r,h_0,h_1}=1\).  The constraints match, for
each grid point \(z_j\), the binned observed probabilities of
\((Y,D)\mid Z=z_j\); match \(\Delta p_j=p(z_j)-p(z_{j-1})\) and
\(\Delta m_j=m(z_j)-m(z_{j-1})\); and impose
\[
\sum_{r,h_0,h_1}\lambda_{r,h_0,h_1}
\one\{r_{j-1}=1,r_j=0\}
\le \int_{z_{j-1}}^{z_j}b(z)\,dz
\]
on each cell.  The objective is the ratio of the weighted positive-transition
treatment-effect sum to the weighted positive-transition mass.  The finite
program has \(2^{J+1}H^2\) nonnegative weights plus one Charnes--Cooper scale
variable and linear constraints after the standard fractional-program
transformation.  As the partition mesh and outcome-bin diameter go to zero,
these finite values converge to \(\theta_L(P,b)\) and \(\theta_U(P,b)\).

\paragraph{Interpretation for Conjecture~\ref{conj:strict-separation}.}
The example isolates the reason the signed derivative ratio is not the
positive-responder target.  On both \(A\) and \(B\), positive responders have
treatment-effect numerator \(a_+/c_+=4/5\).  On \(B\), however, the observed
outcome derivative subtracts the negative-responder numerator
\(a_-=9/250\), while the observed first-stage derivative subtracts the
negative-responder density \(c_-=1/25\).  The observable ratio
\[
q_{\mathrm{IV}}=\frac{\int m'}{\int p'}
\]
therefore averages a signed mixture, not the positive-responder numerator over
positive-responder mass.  The hidden zero-effect cancelling movers explain why
the robustness radius is finite: enlarging the envelope eventually permits
enough invisible opposite movement to dilute the positive-responder ratio down
to \(q_{\mathrm{IV}}\), but the baseline envelope does not permit enough of it.

\paragraph{Synthesis.}
The pointwise inequalities are useful because they reduce every feasible law
to the two observable derivatives \(p'\), \(m'\), the unknown negative response
density \(f=c_-\), and the positive-responder numerator \(u=a_+\).  The path
and level constraints are equally essential: dropping them would give only a
local outer relaxation, not the sharp set; the observable ratio averages a
signed mixture, not the positive-responder numerator over positive-responder
mass.

\section{Step-by-step proof}
\subsection*{Supporting lemmas}
\begin{theorem}[Theorem 1: local signed-derivative envelope]\label{thm:local-envelope}
Under Assumptions~\ref{ass:continuous-iv}--\ref{ass:envelope}, any feasible full-data law must admit an auxiliary negative response density \(f(z)=c_-(z)\) with
\[
\max\{0,-p'(z)\}\le f(z)\le b(z)
\]
almost everywhere, and \(c_+(z)=p'(z)+f(z)\).  Conditional on any such \(f\), the positive-responder numerator density must satisfy
\[
\max\{-p'(z)-f(z),m'(z)-f(z)\}\le a_+(z)\le
\min\{p'(z)+f(z),m'(z)+f(z)\}.
\]
Equivalently, \(a_-(z)=a_+(z)-m'(z)\) lies in \([-f(z),f(z)]\) while \(a_+(z)\) lies in \([-(p'(z)+f(z)),p'(z)+f(z)]\).
\end{theorem}
\begin{proof}
Fix any feasible full-data law and set \(f=c_-\).  Assumption~\ref{ass:envelope} gives \(0\le f\le b\) almost everywhere.  Assumption~\ref{ass:signed-response} gives \(p'=c_+-c_-\), hence
\[
c_+(z)=p'(z)+f(z)
\]
almost everywhere.  Since \(c_+\ge0\), \(p'(z)+f(z)\ge0\), so \(f(z)\ge -p'(z)\).  Combining this with \(f\ge0\) gives
\[
\max\{0,-p'(z)\}\le f(z)\le b(z)
\]
almost everywhere.

Now write \(u=a_+\).  The bounded-outcome numerator restriction in Assumption~\ref{ass:signed-response} gives
\[
-(p'(z)+f(z))\le u(z)\le p'(z)+f(z),
\]
because \(c_+=p'+f\).  The derivative identity \(m'=a_+-a_-\) gives \(a_-=u-m'\).  The negative-response numerator bound \(-c_-\le a_-\le c_-\), with \(c_-=f\), is therefore equivalent to
\[
-f(z)\le u(z)-m'(z)\le f(z),
\]
or
\[
m'(z)-f(z)\le u(z)\le m'(z)+f(z).
\]
Intersecting the interval implied by \(|a_+|\le c_+\) with the interval implied by \(|a_-|\le c_-\) yields
\[
\max\{-p'(z)-f(z),m'(z)-f(z)\}\le u(z)\le
\min\{p'(z)+f(z),m'(z)+f(z)\}.
\]
Substituting back \(u=a_+\) proves the displayed positive-responder numerator constraint.  The equivalent formulation follows from the two separate bounds \(a_-=a_+-m'\in[-f,f]\) and \(a_+\in[-(p'+f),p'+f]\).
\end{proof}

\begin{theorem}[Theorem 2: zero-envelope Heckman--Vytlacil reduction]\label{thm:zero-envelope}
Under Assumptions~\ref{ass:continuous-iv}--\ref{ass:positive-mass}, if \(b(z)=0\) almost everywhere and \(p'(z)\ge\kappa>0\), then \(c_-(z)=0\), \(c_+(z)=p'(z)\), \(a_+(z)=m'(z)\), and
\[
\Theta_{\mathrm{CIV}}(P,0)=
\left\{
\frac{\int_0^1 m'(z)\,dz}{\int_0^1 p'(z)\,dz}
\right\}.
\]
Pointwise, the local effect is the Heckman--Vytlacil LIV/MTE derivative \(m'(z)/p'(z)\) wherever \(p'(z)>0\).
\end{theorem}
\begin{proof}
Let a feasible full-data law be given.  Since \(b=0\) almost everywhere, Assumption~\ref{ass:envelope} implies
\[
0\le c_-(z)\le0
\]
almost everywhere, so \(c_-(z)=0\) almost everywhere.  The first signed-response identity then gives
\[
c_+(z)=p'(z)+c_-(z)=p'(z)
\]
almost everywhere.  The negative numerator bound gives \(-c_-(z)\le a_-(z)\le c_-(z)\), hence \(a_-(z)=0\) almost everywhere.  The outcome derivative identity \(m'=a_+-a_-\) therefore gives \(a_+(z)=m'(z)\) almost everywhere.

By the maintained positive-mass condition and the extra lower bound \(p'(z)\ge\kappa>0\), the denominator is strictly positive:
\[
\int_0^1 c_+(z)\,dz=\int_0^1 p'(z)\,dz\ge\kappa>0.
\]
Substituting the pointwise identities into the target ratio gives, for every feasible law,
\[
\theta(L)=\frac{\int_0^1 a_+(z)\,dz}{\int_0^1 c_+(z)\,dz}
=\frac{\int_0^1 m'(z)\,dz}{\int_0^1 p'(z)\,dz}.
\]
Thus all feasible full-data laws generate the same scalar, so \(\Theta_{\mathrm{CIV}}(P,0)\) is the displayed singleton.  At every margin where \(p'(z)>0\), the same identities give the local ratio \(a_+(z)/c_+(z)=m'(z)/p'(z)\), which is the Heckman--Vytlacil local-IV/MTE derivative in the zero-envelope monotone corner.
\end{proof}

\subsubsection*{Proof of Conjecture 1}
\begin{proof}
Let \(\mathfrak L^+(P,b)\) denote the full-data laws in \(\mathfrak L(P,b)\)
with positive positive-responder mass.  This is the domain already used in
the definition of \(\Theta_{\mathrm{CIV}}(P,b)\) and in the ratio \(R_P\).

First prove necessity.  Fix \(L\in\mathfrak L^+(P,b)\).  Set
\[
f(z)=c_-(z),\qquad u(z)=a_+(z).
\]
By Theorem~\ref{thm:local-envelope}, which uses
Assumptions~\ref{ass:continuous-iv}--\ref{ass:envelope},
\[
\max\{0,-p'(z)\}\le f(z)\le b(z)
\]
and
\[
\max\{-p'(z)-f(z),m'(z)-f(z)\}\le u(z)\le
\min\{p'(z)+f(z),m'(z)+f(z)\}
\]
almost everywhere.  Assumption~\ref{ass:signed-response} gives
\[
c_+(z)=p'(z)+f(z),
\qquad
a_-(z)=u(z)-m'(z),
\]
and Assumption~\ref{ass:path-compatible} gives the full level matching
constraint
\[
\mathcal L_L\bigl(Y(D(z)),D(z)\bigr)=P_{Y,D\mid Z=z}
\quad\text{for every }z\in[0,1],
\]
together with generation of the signed first-stage derivative measure
\(p'(z)\,dz\) and the total-variation response measure
\[
(c_++c_-)\,dz=(p'+2f)\,dz
\]
by the same bounded-variation binary response path.  Therefore
\((f,u)\in\mathfrak S(P,b)\).  Since \(L\in\mathfrak L^+(P,b)\),
\[
\theta(L)=
\frac{\int_0^1 a_+(z)\,dz}{\int_0^1 c_+(z)\,dz}
=
\frac{\int_0^1 u(z)\,dz}
       {\int_0^1(p'(z)+f(z))\,dz}
=R_P(f,u).
\]
It follows that every value in \(\Theta_{\mathrm{CIV}}(P,b)\) is generated by
the displayed optimization.

Now prove sufficiency.  Take any \((f,u)\in\mathfrak S(P,b)\).  By the
definition of the global shadow set, there is a full-data law
\(L\in\mathfrak L^+(P,b)\) whose induced densities satisfy
\[
c_-=f,\qquad a_+=u.
\]
Define
\[
c_+=p'+f,\qquad a_-=u-m'.
\]
The pointwise constraints in Theorem~\ref{thm:local-envelope} give
\[
c_+\ge0,\qquad -c_+\le a_+\le c_+,\qquad -c_-\le a_-\le c_-,
\]
and the two derivative identities are then algebraic:
\[
c_+-c_-=p',\qquad a_+-a_-=m'.
\]
The level constraints in the statement ensure that this same law reproduces
\(P_{Y,D\mid Z=z}\), and the path constraint ensures that the two variation
density measures are produced by one binary bounded-variation response path.
Thus the law is feasible for the original sharp partial-identification
problem, and
\[
R_P(f,u)=\theta(L)\in\Theta_{\mathrm{CIV}}(P,b).
\]
The feasible values of the optimization and the feasible values of the causal
target therefore coincide.

It remains to show that the coincident set is a closed interval and that its
endpoints are attained at extreme laws.  The feasible law set is convex:
mixing two full-data laws that both match \(P\) again matches \(P\), preserves
instrument independence and exclusion, preserves the envelope
\(0\le c_-\le b\), and mixes the response and numerator measures linearly.
For \(L_0,L_1\in\mathfrak L^+(P,b)\), write
\[
N_i=\int_0^1 a_{+,i}(z)\,dz,\qquad
D_i=\int_0^1 c_{+,i}(z)\,dz>0,\qquad
r_i=N_i/D_i .
\]
For \(L_\lambda=\lambda L_0+(1-\lambda)L_1\),
\[
\theta(L_\lambda)=
\frac{\lambda N_0+(1-\lambda)N_1}
     {\lambda D_0+(1-\lambda)D_1}
=
\omega_\lambda r_0+(1-\omega_\lambda)r_1,
\]
where
\[
\omega_\lambda=
\frac{\lambda D_0}{\lambda D_0+(1-\lambda)D_1}.
\]
As \(\lambda\) varies in \([0,1]\), \(\omega_\lambda\) varies continuously
from \(0\) to \(1\).  Hence the image of the convex feasible set under the
positive-denominator ratio contains the whole interval between any two of its
values.

Compactness supplies attainment.  By Assumptions~\ref{ass:differentiable}
and~\ref{ass:envelope},
\[
\int_0^1(c_+(z)+c_-(z))\,dz
\le
\int_0^1 |p'(z)|\,dz+2\int_0^1b(z)\,dz<\infty
\]
uniformly over feasible laws.  Binary bounded-variation paths with variation
bounded by \(M\) are compact in \(L^1([0,1])\) along subsequences by Helly's
selection theorem, and the displayed uniform first-moment bound makes the
family of path laws tight by Markov's inequality.  The outcome space
\([0,1]^2\) is compact by Assumption~\ref{ass:continuous-iv}.  The constraints
of IV validity, exclusion, level matching to \(P_{Y,D\mid Z=z}\), the
absolutely continuous signed derivative identities, and the envelope
constraint are closed under this weak convergence because the induced response
and numerator measures are dominated by
\(|p'|\,dz+2b(z)\,dz\).  Therefore \(\mathfrak L^+(P,b)\) is compact in the
topology generated by bounded continuous path/outcome functionals and the
four aggregate signed-response measures.  The positive-mass convention in
Section~7 makes the denominator functional \(D(L)=\int c_+\,dz\) strictly
positive on this compact feasible set, so it is bounded away from zero.
Consequently \(L\mapsto \theta(L)=N(L)/D(L)\) is continuous and attains its
minimum and maximum.

Let \(K=\mathfrak L^+(P,b)\), \(D(L)=\int c_+\,dz\), and
\(N(L)=\int a_+\,dz\).  If \(\theta_U=\max_K N/D\), then
\[
F_U=\{L\in K:N(L)-\theta_U D(L)=0\}
\]
is a nonempty compact exposed face of \(K\); similarly
\[
F_L=\{L\in K:N(L)-\theta_L D(L)=0\}
\]
is the lower exposed face.  By the Krein--Milman extreme-point theorem, each
nonempty compact convex face contains an extreme point, and an extreme point
of an exposed face is an extreme point of \(K\).  Thus both endpoints are
attained by extreme full-data signed-response laws.  Combining necessity,
sufficiency, interval convexity, and endpoint attainment proves the theorem.
\end{proof}

\subsubsection*{Proof of Conjecture 2}
All response paths in the center law are generated by threshold types whose
thresholds are uniformly distributed over the relevant region.  The visible
positive responders on \(A\) and \(B\) have outcome tags
\[
(Y(0),Y(1))=(1/10,9/10),
\]
so their treatment-effect numerator is \(4/5\) times their response density.
The visible negative responders on \(B\) have outcome tags
\[
(Y(0),Y(1))=(0,9/10),
\]
so their numerator is \(9/10\) times their response density.  The hidden
cancelling responders on \(B\) have
\[
Y(0)=Y(1)=1/2,
\]
and equal positive and negative threshold densities.  The tags
\(0,1/10,1/2,9/10\) are disjoint observed outcome cells, so the derivatives of
the corresponding conditional cell masses identify the visible flows in the
table.  The equal zero-effect cancelling pair at the common tag \(1/2\) moves
the same amount of mass from \(D=0\) to \(D=1\) as from \(D=1\) to \(D=0\);
therefore it leaves the full conditional law \(P_{Y,D\mid Z=z}\) unchanged.
Adding fixed always-treated and never-treated background types with disjoint
tags and sufficiently large slack makes all conditional probabilities strictly
positive and keeps \(p(z)\in(0,1)\) for every \(z\).  Assumptions
\ref{ass:continuous-iv}--\ref{ass:path-compatible} follow directly from this
single bounded-variation threshold construction, IV independence, and bounded
potential outcomes.

By Assumption~\ref{ass:signed-response}, the derivative identities are
\[
p'=c_+-c_-,
\qquad
m'=a_+-a_-.
\]
On \(A\), this gives \(p'=1/10\) and \(m'=2/25\).  On \(B\), the visible
flows give
\[
p'=1/10-1/25=3/50,
\qquad
m'=2/25-9/250=11/250.
\]
Thus \(p'\ge3/50\) on both regions, so the positive first-stage denominator is
bounded away from zero:
\[
\int_0^1p'(z)\,dz
=
\frac12\cdot\frac1{10}+\frac12\cdot\frac3{50}
=\frac2{25}>0.
\]
Substituting the same derivatives into the observable quasi-monotone formula
gives
\[
q_{\mathrm{IV}}(P^\star)
=
\frac{\frac12(2/25)+\frac12(11/250)}
     {\frac12(1/10)+\frac12(3/50)}
=\frac{31}{40}.
\]

Next compute the sharp projection.  Full-\(P\) level matching identifies the
visible tagged flows, as just noted.  The only remaining full-data freedom is
the hidden cancelling density \(s\) on \(B\).  It contributes \(s\) to \(c_+\)
and \(s\) to \(c_-\) on \(B\), contributes no numerator, and is feasible
exactly when the envelope permits the total negative density:
\[
1/25+s\le b_0|_B=21/500.
\]
Hence \(0\le s\le1/500\).  The positive-responder numerator is fixed at
\[
\int_0^1a_+(z)\,dz
=
\frac12\cdot\frac2{25}+\frac12\cdot\frac2{25}
=\frac2{25},
\]
while the positive-responder mass is
\[
\int_0^1c_+(z)\,dz
=
\frac12\cdot\frac1{10}
+\frac12\cdot\left(\frac1{10}+s\right)
=
\frac1{10}+\frac{s}{2}.
\]
Therefore every feasible full-data law has
\[
\theta(s)=
\frac{2/25}{1/10+s/2},
\qquad 0\le s\le1/500.
\]
This continuous one-dimensional image is the interval
\[
\Theta_{\mathrm{CIV}}(P^\star,b_0^\star)
=
\left[
\theta(1/500),\theta(0)
\right]
=
\left[
\frac{80}{101},\frac45
\right].
\]
The distance calculation in the theorem then gives the strict baseline
separation
\[
\operatorname{dist}\!\left(q_{\mathrm{IV}}(P^\star),
\Theta_{\mathrm{CIV}}(P^\star,b_0^\star)\right)
=\frac{69}{4040}>\frac1{100}.
\]

The local-envelope inequalities are slack on the active nonmonotone region
\(B\).  There \(f=c_-=1/25\), \(b_0=21/500\), and \(p'=3/50\), so
\[
0<f<b_0.
\]
For the numerator,
\[
\max\{-p'-f,m'-f\}
=
\max\{-1/10,1/250\}
=\frac1{250}
<
\frac2{25}=a_+,
\]
and
\[
a_+=\frac2{25}
<
\min\{p'+f,m'+f\}
=
\min\{1/10,21/250\}
=\frac{21}{250}.
\]
Thus the strictness claim is compatible with Theorem~\ref{thm:local-envelope}.
On \(A\), \(b_0=0\) by design, so the zero-envelope constraints bind there;
that is the monotone positive-responder part of the witness, not the source of
the separation.

For the breakdown calculation, scaling the envelope to \(t b_0\) changes only
the upper bound on \(s\):
\[
0\le s\le t(21/500)-1/25.
\]
Solving \(\theta(s)=q_{\mathrm{IV}}(P^\star)\) gives
\[
\frac{2/25}{1/10+s/2}=\frac{31}{40}
\quad\Longleftrightarrow\quad
s=\frac1{155}.
\]
Thus
\[
\tau^\star(P^\star;b_0^\star)
=
\frac{1/25+1/155}{21/500}
=
\frac{240}{217}>1.
\]

It remains to record why the property is open.  The finite tagged-cell LP has
one free scalar \(s\), denominator at least \(1/10\), baseline separation
margin \(69/4040\), envelope slack \(1/500\), and breakdown margin
\(240/217-1\).  Perturbing the observable conditional law in integrated total
variation and perturbing \(b_0\) in \(L^1\) perturbs the tagged cell-flow
right-hand sides and the envelope capacity by at most a constant multiple of
the perturbation radius.  Linear-fractional objectives with denominators
bounded below are locally Lipschitz in these right-hand sides; in this
one-dimensional LP this is immediate from the displayed formula for
\(\theta(s)\).  Choose \(\eta>0\) so small that all induced perturbations move
\(q_{\mathrm{IV}}\), the lower endpoint of the baseline interval, and
\(\tau^\star\) by less than one fourth of their displayed strict margins.
Then every \((\widetilde P,\widetilde b_0)\) in the \(L^1\)-ball of radius
\(\eta\) satisfies
\[
\operatorname{dist}\!\left(q_{\mathrm{IV}}(\widetilde P),
\Theta_{\mathrm{CIV}}(\widetilde P,\widetilde b_0)\right)
\ge\frac1{200}
=\delta/2
\]
with \(\delta=1/100\), and its breakdown radius remains strictly larger than
one and locally Lipschitz.  This proves the conjecture.

\section{Sanity checks}
\paragraph{Sanity check for Conjecture~\ref{conj:sharp-pasting}: zero-envelope monotone corner.}
Set \(b=0\) and suppose \(p'(z)\ge\kappa>0\) almost everywhere.  The envelope
forces \(f=c_-=0\).  The local constraints reduce to
\[
c_+(z)=p'(z),\qquad a_-(z)=0,\qquad a_+(z)=m'(z),
\]
which is exactly Theorem~\ref{thm:zero-envelope}.  Therefore
\[
\theta_L(P,0)=\theta_U(P,0)
=
\frac{\int_0^1m'(z)\,dz}{\int_0^1p'(z)\,dz},
\]
and pointwise the local positive-responder effect is the usual
Heckman--Vytlacil derivative \(m'(z)/p'(z)\) wherever \(p'(z)>0\).  Thus the
sharp-pasting theorem collapses to the standard local-IV/MTE formula in the
monotone zero-defier limit.

\paragraph{Sanity check for Conjecture~\ref{conj:strict-separation}.}
If the nonmonotone region is switched off by setting \(b_0=0\) and removing
the visible negative responders on \(B\), then \(s=0\), \(c_-=a_-=0\), and
Theorem~\ref{thm:zero-envelope} applies.  In that corner
\[
p'(z)=1/10,\qquad m'(z)=2/25\quad\text{on both }A\text{ and }B,
\]
so
\[
q_{\mathrm{IV}}(P)=
\frac{\int_0^1m'(z)\,dz}{\int_0^1p'(z)\,dz}
=
\frac{2/25}{1/10}
=\frac45.
\]
The sharp set is the singleton \(\Theta_{\mathrm{CIV}}(P,0)=\{4/5\}\), matching
the zero-envelope Heckman--Vytlacil reduction.  Thus the strict separation is
not an artifact of the two-region partition; it appears only when the negative
response component on \(B\) is genuinely present.

\section{Counterexample or minimality check}
\paragraph{Refutation attempt before proof for Conjecture~\ref{conj:sharp-pasting}.}
The prescribed smallest search is the two-point response-type linear program.
Restrict \(Z\) to two adjacent grid values \(z_0<z_1\) and write
\[
\Delta p=p(z_1)-p(z_0),\qquad
\Delta m=m(z_1)-m(z_0),\qquad
B=\int_{z_0}^{z_1}b(z)\,dz .
\]
The four binary response paths are never-taker \(00\), positive responder
\(01\), negative responder \(10\), and always-taker \(11\).  Let their masses
be \(\pi_{00},\pi_{01},\pi_{10},\pi_{11}\), and let \(A_{01}\) and \(A_{10}\)
be the signed treatment-effect numerators carried by the positive and
negative transition types.  The finite LP constraints are
\[
\Delta p=\pi_{01}-\pi_{10},\qquad
\Delta m=A_{01}-A_{10},\qquad
0\le\pi_{10}\le B,
\]
\[
-\pi_{01}\le A_{01}\le\pi_{01},
\qquad
-\pi_{10}\le A_{10}\le\pi_{10}.
\]
Putting \(f=\pi_{10}\), \(c_+=\pi_{01}=\Delta p+f\), and
\(a_+=A_{01}\), the LP projects to
\[
\max\{0,-\Delta p\}\le f\le B,
\]
\[
\max\{-\Delta p-f,\Delta m-f\}\le a_+\le
\min\{\Delta p+f,\Delta m+f\}.
\]
Thus each candidate vertex of the projected box, namely \(f\) at one of its
two endpoints and \(a_+\) at one of the two induced numerator endpoints,
is represented by a genuine mixture of the four response paths whenever it
satisfies the displayed inequalities.  Conversely, every mixture of the four
paths obeys exactly those inequalities.  The binary search therefore finds no
extreme point outside the conjectured interval and no endpoint that violates
the claimed sharp projection.  The continuous proof below is the compact
bounded-variation analogue of this finite LP calculation; it is the same
Manski--Balke--Pearl extreme-point logic with the adjacent-cell response
types pasted into one path.

\paragraph{Minimality check for Conjecture~\ref{conj:sharp-pasting}: local envelopes alone are not sharp.}
The path-compatibility part of Assumption~\ref{ass:path-compatible} cannot be
dropped.  Let \(p(z)\equiv1\), \(m(z)\equiv1/2\), and take an envelope
\(b(z)\equiv\beta>0\).  The local derivative equations have
\[
p'(z)=0,\qquad m'(z)=0.
\]
The pointwise envelope relaxation alone allows, for every \(z\),
\[
f(z)=\beta,\qquad c_+(z)=\beta,\qquad c_-(z)=\beta,\qquad
a_+(z)=a_-(z)=0.
\]
These choices satisfy the displayed inequalities in
Theorem~\ref{thm:local-envelope}.  They cannot be induced by a full-data law
matching \(P_{D\mid Z=z}\) with \(p(z)=1\) for every \(z\): when \(D(z)=1\)
almost surely at a margin, there is no mass in state \(0\) available to make a
positive \(0\to1\) jump at that same margin.  Hence \(c_+(z)>0\) on a set of
positive measure contradicts generation by one binary response path matching
the level \(p(z)=1\).  This example shows that the local signed-derivative
box is only an outer relaxation unless full-\(P\) level matching and global
path compatibility are imposed.  The theorem's sharpness conclusion therefore
uses exactly the extra global condition that the conjecture names.

\paragraph{Refutation attempt before proof for Conjecture~\ref{conj:strict-separation}.}
The required refutation-first check was run on the smallest two-region
response-type LP.  The first candidate was the derivative-only relaxation from
Theorem~\ref{thm:local-envelope}.  With only the variables \(f=c_-\) and
\(u=a_+\), positive \(p'\), and \(q_{\mathrm{IV}}\in[-1,1]\), the choice
\(f=0\) and \(u=m'\) puts the quasi-monotone formula back into the relaxed
set.  This candidate therefore cannot witness strict separation; it shows that
full-\(P\) level information is essential.

The second candidate was the tagged two-region LP used in
Theorem~\ref{thm:strict-separation}.  Its decision variable is the hidden
cancelling density \(s\) on \(B\), subject at baseline to
\[
0\le s\le1/500.
\]
The LP vertices are \(s=0\) and \(s=1/500\), giving
\[
R(0)=4/5,\qquad R(1/500)=80/101.
\]
The quasi-monotone value is \(31/40\), and
\[
31/40<80/101<4/5.
\]
Thus vertex enumeration produces no extreme point contradicting the conjecture.
It instead produces the positive witness: the conjectured separation holds at
the center, and the unique breakdown vertex occurs only after the envelope is
scaled to \(t=240/217>1\).

\paragraph{Counterexample or minimality check for Conjecture~\ref{conj:strict-separation}.}
Dropping full-\(P\) level matching destroys the result.  In the same center
law, keep only \(p'\), \(m'\), and the envelope.  Since \(p'(z)>0\) and
\(|m'(z)|\le p'(z)\) on both regions, the local relaxation admits the monotone
completion \(f=0\) and \(a_+=m'\), which yields
\[
\frac{\int_0^1a_+(z)\,dz}{\int_0^1(p'(z)+f(z))\,dz}
=
q_{\mathrm{IV}}(P).
\]
Hence the separation is a full-distribution and path-compatibility statement,
not a consequence of the derivative envelope alone.  The support condition
needed in the proof is the separated tagged-cell condition: the outcome cells
used to identify the visible positive and negative response flows must remain
disjoint with positive probability margins.  Without that condition, an
alternative monotone or weakly nonmonotone completion can absorb the signed
derivative ratio into the feasible projection.

\section{Formalization checklist}
\begin{enumerate}[label=\arabic*.]
\item \(\mathfrak L(P,b)\).  Typed signature: observable law \(P\), envelope
\(b\in L^1_+([0,1])\) \(\mapsto\) set of full-data laws over bounded-variation
binary response paths and bounded potential outcomes.  Role: original feasible
set for the causal target.  Dependencies: Assumptions~\ref{ass:continuous-iv}--\ref{ass:path-compatible}.

\item \(\mathfrak S(P,b)\).  Typed signature:
\((P,b)\mapsto\{(f,u)\in L^1_+([0,1])\times L^1([0,1])\}\).  Role: global
shadow set for negative response density and positive-responder numerator.
Dependencies: existence of a full-data law in \(\mathfrak L(P,b)\), full-\(P\)
level matching, path compatibility, and positive-mass convention.

\item \(R_P(f,u)\).  Typed signature:
\(\mathfrak S(P,b)\to\R\), \(R_P(f,u)=\int u/\int(p'+f)\).  Role:
computable ratio optimized for \(\theta_L(P,b)\) and \(\theta_U(P,b)\).
Dependencies: \(p'\in L^1\), \(f\in L^1_+\), and
\(\int_0^1(p'(z)+f(z))\,dz>0\).

\item \(\Theta_{\mathrm{CIV}}(P,b)\).  Typed signature: \((P,b)\mapsto\)
subset of \(\R\).  Role: sharp positive-responder identified set.
Dependencies: \(\mathfrak L(P,b)\), the target map \(\theta(L)\), and the
positive positive-responder mass condition.

\item Local envelope inequalities.  Typed signature:
\((p',m',b)\mapsto\) pointwise interval constraints for \(f=c_-\) and
\(u=a_+\).  Role: necessary local accounting identities and finite-cell LP
projection.  Dependencies: signed response decomposition, bounded outcomes,
and calibrated envelope.

\item Finite mesh linear-fractional program.  Typed signature: partition
\(\pi\), outcome grid \(\mathcal H\), binned observed laws, increments
\((\Delta p_j,\Delta m_j)\), and cell envelope masses \(\mapsto\) lower and
upper objective values.  Role: computable approximation and extreme-point
description of the endpoint formula.  Dependencies: nonnegative type weights
\(\lambda_{r,h_0,h_1}\), observed cell matching constraints, response
increment constraints, envelope constraints, and Charnes--Cooper
linearization.

\item Tagged two-region witness.  Typed signature: regions \(A,B\), table of
constant densities \((c_+,a_+,c_-,a_-)\), envelope \(b_0\), and scalar hidden
cancelling density \(s\).  Role: explicit open-set witness for strict
separation and breakdown radius \(240/217\).  Dependencies: separated outcome
tags, \(0\le s\le1/500\) at baseline, denominator lower bound, and local
Lipschitz perturbation of the one-dimensional linear-fractional program.

\item Bypass-justified substrate audit record.  Typed signature: queried names
\(\{\texttt{POManskiIVSystem},\texttt{manski\_bounds\_ATE},\texttt{BalkePearl}\}\)
\(\mapsto\) no local declaration matches from the available index.  Role:
documents why this Stage 0 note stays at the paper-level algebraic and
measure-path abstraction rather than citing repository primitives.  Dependencies:
no Lean source files were opened or cited.
\end{enumerate}

\section{Conclusion and references}
Conjecture 1 (sharp-pasting) is confirmed; Under Assumptions~\ref{ass:continuous-iv}--\ref{ass:path-compatible}, the sharp identified set \(\Theta_{\mathrm{CIV}}(P,b)\) is exactly the interval obtained by optimizing.

Conjecture 2 (strict-separation) is confirmed; Conjecture 2 is confirmed: the two-region tagged-response witness class below gives a nonempty \(L^1\)-open neighborhood on which \(q_{\mathrm{IV}}\) is separated from \(\Theta_{\mathrm{CIV}}(P,b_0)\) and has locally Lipschitz breakdown radius strictly larger than one.

\begin{thebibliography}{99}
\bibitem{BalkePearl1997}
Balke, A. and Pearl, J. (1997).
Bounds on treatment effects from studies with imperfect compliance.
\emph{Journal of the American Statistical Association}, 92(439), 1171--1176.

\bibitem{DeChaisemartin2017}
de Chaisemartin, C. (2017).
Tolerating defiance? Local average treatment effects without monotonicity.
\emph{Quantitative Economics}, 8(2), 367--396.

\bibitem{HeckmanVytlacil2005}
Heckman, J. J. and Vytlacil, E. (2005).
Structural equations, treatment effects, and econometric policy evaluation.
\emph{Econometrica}, 73(3), 669--738.

\bibitem{ImbensAngrist1994}
Imbens, G. W. and Angrist, J. D. (1994).
Identification and estimation of local average treatment effects.
\emph{Econometrica}, 62(2), 467--475.

\bibitem{KennedyLorchSmall2019}
Kennedy, E. H., Lorch, S., and Small, D. S. (2019).
Robust causal inference with continuous instruments using the local instrumental variable curve.
\emph{Journal of the Royal Statistical Society: Series B}, 81(1), 121--143.

\bibitem{Kitagawa2015}
Kitagawa, T. (2015).
A test for instrument validity.
\emph{Econometrica}, 83(5), 2043--2063.

\bibitem{Klein2010}
Klein, T. J. (2010).
Heterogeneous treatment effects: Instrumental variables without monotonicity?
\emph{Journal of Econometrics}, 155(2), 99--116.

\bibitem{Manski1990}
Manski, C. F. (1990).
Nonparametric bounds on treatment effects.
\emph{American Economic Review Papers and Proceedings}, 80(2), 319--323.

\bibitem{Marx2024}
Marx, P. (2024).
Sharp bounds for marginal treatment effects.
\emph{Working paper}.

\bibitem{MogstadSantosTorgovitsky2018}
Mogstad, M., Santos, A., and Torgovitsky, A. (2018).
Using instrumental variables for inference about policy relevant treatment parameters.
\emph{Econometrica}, 86(5), 1589--1619.

\bibitem{Vytlacil2002}
Vytlacil, E. (2002).
Independence, monotonicity, and latent index models: An equivalence result.
\emph{Econometrica}, 70(1), 331--341.

\bibitem{Yap2022}
Yap, C. (2022).
Sensitivity analysis for monotonicity in instrumental variables.
\emph{Working paper}.
\end{thebibliography}

\end{document}
